% =====================================================================
% Semana 4 — MOSFET de potencia con IRF640N
% Compilar desde docs/presentations con pdflatex (2 pasadas)
% =====================================================================
\documentclass[aspectratio=169]{beamer}

\usepackage[utf8]{inputenc}
\usepackage[T1]{fontenc}
\usepackage{lmodern}
\usepackage[spanish,es-tabla]{babel}
\usepackage{amsmath,amssymb}
\usepackage{siunitx}
\usepackage{booktabs}
\usepackage{microtype}
\usepackage{circuitikz}

\definecolor{navy}{RGB}{23,54,93}
\definecolor{navysoft}{RGB}{72,101,140}
\definecolor{gold}{RGB}{179,142,62}
\definecolor{ink}{RGB}{44,51,58}
\definecolor{soft}{RGB}{244,241,234}
\definecolor{warn}{RGB}{150,44,35}

\usetheme{default}
\setbeamertemplate{navigation symbols}{}
\setbeamertemplate{footline}[frame number]
\setbeamercolor{structure}{fg=navy}
\setbeamercolor{frametitle}{fg=white,bg=navy}
\setbeamercolor{title}{fg=navy}
\setbeamercolor{subtitle}{fg=navysoft}
\setbeamercolor{itemize item}{fg=gold}
\setbeamercolor{itemize subitem}{fg=gold}
\setbeamercolor{block title}{fg=white,bg=navy}
\setbeamercolor{block body}{bg=soft}
\setbeamercolor{alerted text}{fg=warn}
\setbeamertemplate{itemize item}{\scriptsize\raise1.25pt\hbox{\donotcoloroutermaths$\bullet$}}
\setbeamertemplate{itemize subitem}{\scriptsize\raise1.25pt\hbox{\donotcoloroutermaths$\circ$}}
\setbeamerfont{frametitle}{size=\large,series=\bfseries}

\title[MOSFET]{MOSFET de potencia: del datasheet al conmutador}
\subtitle{IRF640N · Semana 4 · Fundamentos de Sistemas Electrónicos Analógicos}
\institute{Licenciatura en Ingeniería Aeroespacial · UNAM}
\date{Semestre 2027-1}

\begin{document}

\begin{frame}[plain]
  \titlepage
\end{frame}

\begin{frame}{Pregunta de entrada}
  \begin{columns}[T,onlytextwidth]
    \begin{column}{0.47\textwidth}
      \begin{block}{Afirmación}
        ``El IRF640N tiene $V_{GS(th)}=2$--$4\,\mathrm{V}$; por tanto un GPIO
        de \SI{3.3}{\volt} lo enciende por completo.''
      \end{block}
    \end{column}
    \begin{column}{0.47\textwidth}
      \textbf{Antes de seguir:}
      \begin{enumerate}
        \item ¿Verdadero o falso?
        \item ¿Qué fila del datasheet usarías para demostrarlo?
        \item ¿Qué variable falta para hablar de pérdidas?
      \end{enumerate}
    \end{column}
  \end{columns}
  \vfill
  \begin{alertblock}{Meta de la clase}
    Separar \textbf{inicio de conducción}, \textbf{drive garantizado} y
    \textbf{límite absoluto}.
  \end{alertblock}
\end{frame}

\begin{frame}{Objetivos observables}
  Al finalizar podrás:
  \begin{itemize}
    \item identificar corte, región óhmica y saturación con sus condiciones;
    \item decidir si una tensión de compuerta garantiza baja $R_{DS(on)}$;
    \item estimar pérdida de conducción y esfuerzo del driver mediante $Q_g$;
    \item dibujar el camino de una corriente inductiva al apagar;
    \item convertir corriente en tensión para un ADC y evaluar su resolución;
    \item explicar qué puede y qué no puede validar un modelo SPICE.
  \end{itemize}
\end{frame}

\begin{frame}{1 · Estructura y control por campo eléctrico}
  \begin{columns}[T]
    \begin{column}{0.52\textwidth}
      \centering
      \begin{tikzpicture}[scale=0.83]
        \draw[fill=gray!18] (0,0) rectangle (5.1,2.5);
        \draw[fill=blue!25] (0.5,0.15) rectangle (1.55,2.1);
        \draw[fill=blue!25] (3.55,0.15) rectangle (4.6,2.1);
        \draw[fill=orange!30] (0.35,1.1) rectangle (4.75,1.5);
        \draw[very thick,fill=white] (0.3,1.55) rectangle (4.8,1.72);
        \draw[very thick] (1.7,1.92) -- (3.4,1.92);
        \draw (2.55,1.92) -- (2.55,2.75) node[above] {G};
        \draw (1.0,2.1) -- (1.0,2.75) node[above] {S};
        \draw (4.1,2.1) -- (4.1,2.75) node[above] {D};
        \node at (2.55,0.52) {sustrato P};
        \node[navy] at (2.55,1.30) {canal inducido};
      \end{tikzpicture}
    \end{column}
    \begin{column}{0.44\textwidth}
      \begin{itemize}
        \item El óxido aísla el gate: $I_G\approx0$ en estado estacionario.
        \item $V_{GS}$ crea y modula el canal.
        \item El gate es una \textbf{carga capacitiva}; en cada flanco el driver
          entrega o retira carga.
      \end{itemize}
      \begin{block}{No olvidar}
        Control por tensión no significa driver sin corriente.
      \end{block}
    \end{column}
  \end{columns}
\end{frame}

\begin{frame}{2 · Regiones del NMOS}
  \small
  \begin{tabular}{p{0.16\textwidth}p{0.28\textwidth}p{0.42\textwidth}}
    \toprule
    Región & Condición ideal & Modelo de primer orden \\
    \midrule
    \textbf{Corte} & $V_{GS}<V_{th}$ & $I_D\approx0$ \\
    \textbf{Óhmica / triodo} & $V_{GS}>V_{th}$ y $0\le V_{DS}<V_{GS}-V_{th}$ &
      $I_D=k_n\!\left[(V_{GS}-V_{th})V_{DS}-\tfrac12V_{DS}^{2}\right]$ \\
    \textbf{Saturación} & $V_{GS}>V_{th}$ y $V_{DS}\ge V_{GS}-V_{th}$ &
      $I_D\approx\tfrac12k_n(V_{GS}-V_{th})^2$ \\
    \bottomrule
  \end{tabular}
  \vfill
  \begin{block}{Convención de esta semana}
    Interruptor \textbf{ON} $\Rightarrow$ región \textbf{óhmica}, con
    $V_{DS}\approx I_D R_{DS(on)}$. Saturación MOSFET no equivale a
    saturación BJT.
  \end{block}
\end{frame}

\begin{frame}{3 · Umbral no es encendido}
  \begin{columns}[T]
    \begin{column}{0.46\textwidth}
      \begin{block}{$V_{GS(th)}=2$--$4\,\mathrm{V}$}
        Se mide con $I_D=\SI{250}{\micro\ampere}$ y $V_{DS}=V_{GS}$.
        Detecta el \textbf{inicio} del canal.
      \end{block}
      \begin{alertblock}{A \SI{3.3}{\volt}}
        El fabricante no garantiza $R_{DS(on)}$ para el IRF640N.
      \end{alertblock}
    \end{column}
    \begin{column}{0.48\textwidth}
      \begin{block}{$R_{DS(on)}\le\SI{0.15}{\ohm}$}
        Se especifica con $V_{GS}=\SI{10}{\volt}$, $I_D=\SI{11}{\ampere}$ y
        $T_J=\SI{25}{\celsius}$.
      \end{block}
      \begin{exampleblock}{Decisión}
        Para esta práctica: driver de 10--12 V. Para drive directo de 3.3 V:
        escoger otro MOSFET con $R_{DS(on)}$ garantizada a 3.3 V.
      \end{exampleblock}
    \end{column}
  \end{columns}
  \vfill
  {\tiny Fuente: datasheet IRF640N de Infineon, características eléctricas.}
\end{frame}

\begin{frame}{4 · Leer el datasheet con condiciones}
  \small
  \begin{tabular}{p{0.26\textwidth}p{0.20\textwidth}p{0.44\textwidth}}
    \toprule
    Dato & Valor & Qué significa \\
    \midrule
    $V_{DS,max}$ & \SI{200}{\volt} & límite absoluto; requiere margen y revisión de SOA \\
    $V_{GS,max}$ & $\pm\SI{20}{\volt}$ & límite gate-source, no tensión recomendada \\
    $I_D$ publicado & \SI{18}{\ampere} & con $T_C=\SI{25}{\celsius}$ y condiciones térmicas \\
    $Q_g$ máximo & \SI{67}{\nano\coulomb} & $I_D=\SI{11}{\ampere}$, $V_{DS}=\SI{160}{\volt}$, $V_{GS}=\SI{10}{\volt}$ \\
    $C_{iss}$ típico & \SI{1160}{\pico\farad} & $V_{GS}=0$, $V_{DS}=\SI{25}{\volt}$, \SI{1}{\mega\hertz} \\
    \bottomrule
  \end{tabular}
  \begin{block}{Regla de lectura}
    Un número sin condición de prueba está incompleto. Un máximo absoluto no
    es un punto normal de operación.
  \end{block}
\end{frame}

\begin{frame}{5 · NMOS de lado bajo}
  \begin{columns}[T,onlytextwidth]
    \begin{column}{0.42\textwidth}
      \centering
      \resizebox{0.90\linewidth}{!}{%
      \begin{circuitikz}[american]
        \node[nmos] (M) at (0,1.3) {};
        \draw (M.S) -- (0,0) node[ground]{};
        \draw (M.D) -- (0,2.55);
        \draw (0,3.0) circle (0.45) node {M};
        \draw (0,3.45) -- (0,3.8);
        \node[above] at (0,3.8) {$+5\,\mathrm{V}$};
        \node[right,align=left] at (0.55,3.0) {motor DC, 6 V nom.\\250--400 mA};
        \draw (M.G) to[R,l_={$R_G=100\,\Omega$}] (-2.3,1.3);
        \node[left,align=right] at (-2.3,1.3) {driver\\0/12 V};
        \draw[->,gold,very thick] (0.45,3.0) -- (0.45,2.2) node[midway,right] {$I_D$};
        \node[right] at (0.5,1.55) {D};
        \node[right] at (0.5,0.75) {S};
      \end{circuitikz}
      }
    \end{column}
    \begin{column}{0.54\textwidth}
      \footnotesize
      \begin{tabular}{@{}lll@{}}
        \toprule
        & OFF & ON \\
        \midrule
        $V_{GS}$ & $0$ & $\approx\SI{12}{\volt}$ \\
        $I_D$ & $\approx0$ & $0.25$--$0.40\,\mathrm{A}$ \\
        $V_D$ & $\approx\SI{5}{\volt}$ & $\le60\,\mathrm{mV}$ \\
        Región & corte & óhmica \\
        \bottomrule
      \end{tabular}
      \vspace{0.6em}
      \begin{block}{Inversión}
        Gate alto $\Rightarrow$ drenaje bajo. Gate bajo $\Rightarrow$ drenaje alto.
      \end{block}
    \end{column}
  \end{columns}
\end{frame}

\begin{frame}{6 · Dos puntos nominales del motor}
  Motor nominal de 6 V, alimentado con 5 V. Como envolvente conservadora,
  $5/0.25=\SI{20}{\ohm}$ sin carga y $5/0.40=\SI{12.5}{\ohm}$ con carga. Con
  $R_{DS(on),max}=\SI{0.15}{\ohm}$ a \SI{25}{\celsius}:
  \[
    I_{0}\approx\frac{5}{20+0.15}=\SI{248.1}{\milli\ampere},\qquad
    I_{carga}\approx\frac{5}{12.5+0.15}=\SI{395.3}{\milli\ampere}
  \]
  \[
    V_{DS(on)}\le\SI{59.3}{\milli\volt},\qquad
    P_{cond}\le\SI{23.4}{\milli\watt}\quad\text{(punto cargado)}
  \]
  \begin{alertblock}{Alcance}
    Si 250/400 mA son de ficha a 6 V, use 24/15 $\Omega$ y mida a 5 V.
    \SI{400}{\milli\ampere} tampoco es la corriente de bloqueo.
  \end{alertblock}
\end{frame}

\begin{frame}{7 · La compuerta es capacitiva}
  \begin{columns}[T]
    \begin{column}{0.48\textwidth}
      \[
        I_{G,avg}\approx Q_g f
        =\SI{67}{\nano\coulomb}\cdot\SI{100}{\kilo\hertz}
        =\SI{6.7}{\milli\ampere}
      \]
      \[
        P_{driver}\approx Q_g V_{drive}f\approx\SI{80}{\milli\watt}
      \]
    \end{column}
    \begin{column}{0.46\textwidth}
      \begin{itemize}
        \item $R_G$ limita corriente pico y amortigua oscilación.
        \item La meseta Miller ocurre mientras cambia $V_{DS}$.
        \item Para tiempos de conmutación se usa la curva de $Q_g$, no solo
          $C_{iss}$.
      \end{itemize}
    \end{column}
  \end{columns}
  \begin{block}{Predicción}
    Si $f$ aumenta diez veces, ¿qué ocurre con la corriente media y la potencia
    del driver?
  \end{block}
\end{frame}

\begin{frame}{8 · Carga inductiva: la corriente busca camino}
  \begin{columns}[T]
    \begin{column}{0.44\textwidth}
      \centering
      \begin{circuitikz}[american]
        \node[nmos] (M) at (0,0.9) {};
        \draw (M.S) -- (0,-0.2) node[ground]{};
        \draw (M.D) -- (0,2.35);
        \draw (0,2.85) circle (0.50) node {M};
        \draw (0,3.35) -- (0,3.7);
        \node[above] at (0,3.7) {$+5\,\mathrm{V}$};
        \draw (0,1.75) -- (2.0,1.75) to[D,l={$D_F$}] (2.0,3.7) -- (0,3.7);
        \draw (M.G) -- (-1.1,0.9) node[left] {driver};
        \draw[->,gold,very thick] (1.55,2.05) -- (1.55,3.2);
      \end{circuitikz}
    \end{column}
    \begin{column}{0.52\textwidth}
      \begin{itemize}
        \item Diodo: ánodo al drenaje, cátodo a $+5\,\mathrm{V}$.
        \item Durante ON está inversamente polarizado.
        \item Al apagar, conduce el mismo sentido de corriente de la bobina.
        \item El drenaje queda cerca de $V_{CC}+V_F(I,T)$.
      \end{itemize}
      \[
        E_L=\tfrac12L_{motor}I^2
      \]
      El modelo didáctico usa $L_{place}=\SI{100}{\milli\henry}$ y
      $R_{eq}=\SI{12.5}{\ohm}$: $\tau\approx\SI{7.9}{\milli\second}$ y
      $E\approx\SI{7.7}{\milli\joule}$. Se debe medir la $L$ real.
    \end{column}
  \end{columns}
\end{frame}

\begin{frame}{9 · Con diodo y sin diodo}
  \begin{columns}[T]
    \begin{column}{0.47\textwidth}
      \begin{exampleblock}{Con flyback}
        La energía se disipa en la resistencia de la bobina y el diodo. El
        drenaje queda cerca de $\SI{5}{\volt}+V_F(I,T)$. El netlist da
        $V_{DS,max}=\SI{5.901}{\volt}$ e $I_{fly,pico}=\SI{393.6}{\milli\ampere}$.
      \end{exampleblock}
    \end{column}
    \begin{column}{0.47\textwidth}
      \begin{alertblock}{Sin flyback}
        $V_{DS}$ aumenta hasta que una capacitancia, un clamp, la avalancha o
        una falla ofrece camino. No existe un pico universal.
      \end{alertblock}
    \end{column}
  \end{columns}
  \vfill
  \textbf{Interpretación correcta:} el archivo sin diodo usa un clamp
  pedagógico cerca de \SI{200}{\volt} para mantener una solución numérica; no
  predice supervivencia ni energía de avalancha.
\end{frame}

\begin{frame}{10 · De corriente a tensión para el ADC}
  \begin{columns}[T]
    \begin{column}{0.46\textwidth}
      \resizebox{0.94\linewidth}{!}{%
      \begin{circuitikz}[american]
        \node[nmos] (M) at (0,1.7) {};
        \draw (M.D) -- (0,3.15);
        \draw (0,3.55) circle (0.40) node {M};
        \draw (0,3.95) -- (0,4.1);
        \node[above] at (0,4.1) {$+5\,\mathrm{V}$};
        \draw (M.S) -- (0,0.9) to[R,l={$R_S=1\,\Omega$}] (0,-0.5) node[ground]{};
        \draw (0,0.9) -- (2.0,0.9) to[R,l={$10\,\mathrm{k}\Omega$}] (3.7,0.9)
          to[R,l={$10\,\mathrm{k}\Omega$}] (3.7,-0.5) node[ground]{};
        \node[above] at (3.7,0.9) {ADC};
      \end{circuitikz}
      }
    \end{column}
    \begin{column}{0.50\textwidth}
      \footnotesize
      Con $R_S=\SI{1}{\ohm}$, el shunt cambia el punto de operación:
      \begin{center}
      \begin{tabular}{@{}lcc@{}}
        \toprule
        & sin carga & con carga \\
        \midrule
        $I_D$ & 236.4 mA & 366.3 mA \\
        $V_S$ & 236.4 mV & 366.3 mV \\
        $V_{ADC}$ & 118.2 mV & 183.2 mV \\
        \bottomrule
      \end{tabular}
      \end{center}
      \begin{block}{Resolución}
        ADC ideal de 12 bits y 3.3 V: 147/227 cuentas. El shunt disipa
        aproximadamente 55.9/134.2 mW.
      \end{block}
    \end{column}
  \end{columns}
\end{frame}

\begin{frame}{11 · Medir bien: $V_{GS}$ no siempre es $V_G$}
  \begin{itemize}
    \item Con un shunt de fuente, $V_S>0$ y $V_{GS}=V_G-V_S$.
    \item La simulación mide $V(\text{gate},\text{sense})$ y
      $V(\text{drain},\text{sense})$.
    \item El divisor 10 k$\Omega$/10 k$\Omega$ atenúa por 0.5; no puede dar
      ganancia mayor que uno.
    \item Protección ante fallas requiere análisis adicional: resistencia de
      entrada, clamps, rango común y corriente de inyección del ADC.
  \end{itemize}
  \begin{alertblock}{Pregunta}
    ¿Qué error aparece si se reporta $V_G=12\,\mathrm{V}$ como si fuera
    $V_{GS}$ cuando $V_S\approx\SI{367}{\milli\volt}$?
  \end{alertblock}
\end{frame}

\begin{frame}{12 · Qué valida cada modelo}
  \small
  \begin{tabular}{p{0.24\textwidth}p{0.28\textwidth}p{0.34\textwidth}}
    \toprule
    Herramienta & Sí ayuda a estudiar & No autoriza por sí sola \\
    \midrule
    Cálculo con datasheet & peor caso de conducción y límites & forma de onda dinámica completa \\
    Sustituto pedagógico & nodos, polaridades, causalidad, medidas & SOA, térmica, avalancha, EMI \\
    Modelo validado & comportamiento dentro de su alcance documentado & ignorar tolerancias o montaje \\
    Prototipo medido & parásitos y hardware real & exceder límites o eliminar protecciones \\
    \bottomrule
  \end{tabular}
  \begin{block}{Nuestro archivo}
    \texttt{semana4\_irf640n\_educativo.lib} es deliberadamente un sustituto
    de aprendizaje, no un macromodelo oficial de Infineon.
  \end{block}
\end{frame}

\begin{frame}{13 · MOSFET y BJT: comparar sin absolutos}
  \small
  \begin{tabular}{p{0.23\textwidth}p{0.32\textwidth}p{0.35\textwidth}}
    \toprule
    Criterio & BJT NPN & NMOS \\
    \midrule
    Entrada & corriente de base & carga/descarga de gate \\
    ON & $V_{CE(sat)}$ a un $I_C/I_B$ dado & $I_D R_{DS(on)}$ a un $V_{GS}$ dado \\
    Dinámica & carga almacenada si satura profundo & $Q_g$ y meseta Miller \\
    Diseño & verificar base, ganancia forzada y disipación & verificar drive, SOA, térmica y transitorios \\
    Usos & ambos pueden amplificar o conmutar & depende de tensión, corriente y frecuencia \\
    \bottomrule
  \end{tabular}
\end{frame}

\begin{frame}{14 · Flujo de trabajo y herramientas}
  \begin{enumerate}
    \item Definir $V$, $I$, carga, frecuencia, temperatura y fallas.
    \item Leer límites, SOA, $R_{DS(on)}$ a la tensión real y $Q_g$.
    \item Calcular conducción, driver, protección y térmica.
    \item Simular primero cada bloque en LTspice.
    \item Cuando exista, usar un modelo validado del fabricante; InfineonSpice
      ofrece una biblioteca de componentes.
    \item Construir con límites de corriente y medir con masa común.
  \end{enumerate}
  \begin{block}{Orden sano}
    Predice $\rightarrow$ simula $\rightarrow$ explica la diferencia
    $\rightarrow$ decide si el hardware es seguro.
  \end{block}
\end{frame}

\begin{frame}{15 · Contexto aeroespacial y térmico}
  \begin{itemize}
    \item Calentadores, solenoides, actuadores y buses auxiliares combinan
      cargas resistivas, inductivas y capacitivas.
    \item Seleccionar con margen de tensión, SOA, transitorios, radiación,
      temperatura y modo de falla; no solo con corriente nominal.
    \item En vacío no hay convección, pero sí conducción al montaje y radiación.
      Un $R_{\theta JA}$ de aire quieto no se transfiere automáticamente.
    \item El diodo flyback puede ralentizar la liberación de un actuador; una
      red de clamp distinta puede ser necesaria si el tiempo importa.
  \end{itemize}
\end{frame}

\begin{frame}{16 · Ruta de laboratorio}
  \begin{enumerate}
    \item PWM resistivo y carga equivalente de gate.
    \item Bobina con flyback y contraste numérico sin diodo.
    \item Shunt de fuente y escalado al ADC.
    \item Conmutador BJT corregido con dos constantes de tiempo.
  \end{enumerate}
  \begin{block}{Evidencia por etapa}
    Predicción, captura de ondas, tabla \texttt{.meas}, diferencia contra el
    cálculo, limitación del modelo y decisión de hardware.
  \end{block}
  \vfill
  \centering
  \texttt{docs/labs/semana\_4\_conmutador\_potencia\_mosfet.md}
\end{frame}

\begin{frame}{Recuperación final}
  Responde sin apuntes:
  \begin{enumerate}
    \item ¿Por qué $V_{GS(th)}$ no decide si un GPIO de 3.3 V es suficiente?
    \item ¿En qué región está un NMOS usado como switch ON?
    \item ¿Por qué $Q_g$ importa si $I_G\approx0$ en DC?
    \item ¿Qué camino sigue la corriente de la bobina al apagar?
    \item ¿Qué sacrifica un divisor que lleva aproximadamente
      \SI{237}{\milli\volt}--\SI{367}{\milli\volt} a la mitad?
    \item ¿Qué resultado SPICE no aceptarías sin revisar el alcance del modelo?
  \end{enumerate}
\end{frame}

\begin{frame}{Fuentes y material de apoyo}
  \begin{itemize}
    \item \href{https://www.infineon.com/assets/row/public/documents/24/49/infineon-irf640n-datasheet-en.pdf}{Infineon: IRF640N datasheet}.
    \item \href{https://www.infineon.com/design-resources/simulation-modeling/infineonspice-offline-simulation-tool}{InfineonSpice: simulador y biblioteca de modelos}.
    \item \href{https://www.analog.com/en/resources/design-tools-and-calculators/ltspice-simulator.html}{Analog Devices: LTspice}.
    \item Guía, netlists verificados, hoja activa y pizarrones de la Semana 4
      incluidos en el repositorio del curso.
  \end{itemize}
  \begin{alertblock}{Criterio final}
    Para límites eléctricos manda el datasheet; para aprender causalidad ayuda
    el modelo; para aprobar hardware se necesitan ambos más verificación.
  \end{alertblock}
\end{frame}

\end{document}
